% Short running form; the full title is too wide for the head's right box.
\runninghead{Thanks at Both Ends}
\properlane{nothing-of-our-own}{Thanks at Both Ends, and a Man Who Was Carried: The Grace the Collect Confesses}
\label{sec:nothing-of-our-own}

Heard from its Collect, the Mass is unusually plain about the grammar of grace,
and it is plain in its structure as well as in its words. The lesson's first
clause is a thanksgiving for a grace already given: \latin{Gratias ago Deo meo
semper pro vobis in gratia Dei, quae data est vobis}. The last prayer of the
Mass is a thanksgiving for a gift just received: \latin{Gratias tibi referimus,
Domine, sacro munere vegetati}. The single oration between them states the
reason for both: \latin{quia tibi sine te placere non possumus}. And the man at
the centre of the Gospel is carried in, spoken to, forgiven and healed without
saying a word. The petitions of such a Mass are all of one kind: that God
should do in us what we cannot begin. St John Chrysostom, reading the Epistle
clause by clause, takes the Corinthians' gifts out of their hands and gives
them back to God; St Augustine, reading the Communion, names the one present a
creature may bring; and St Jerome and Chrysostom, reading the Gospel, disagree
about whose faith brought the paralytic to Christ, a disagreement this reading
keeps.

\subsection{Thanks for a grace already given}

St Paul, Chrysostom observes, begins almost every letter with thanksgiving, and
here the occasion is more urgent than elsewhere: ``For he that gives thanks,
does so, both as being well off, and as in acknowledgment of a favor: now a
favor is not a debt nor a requital nor a payment: which indeed every where is
important to be said, but much more in the case of the Corinthians who were
gaping after the dividers of the Church.'' Each phrase of the lesson then
becomes a correction. On ``the grace of God'': ``Do you see how from every
quarter he draws topics for correcting them? For where grace is, works are not;
where works, it is no more grace. If therefore it be grace, why are you
high-minded? Whence is it that you are puffed up?'' On ``which is given you'':
``And by whom was it given? By me, or by another Apostle? Not at all, but by
Jesus Christ.'' And of the praises that follow he says that they are only
praises in appearance: ``I say, seem; for not even did this praise belong to
them, but to the grace of God. For that they had remission of sins, and were
justified, this was of the Gift from above'' (\work{Homilies on First
Corinthians} 2, 1--5). In the Epistle, then, what the Corinthians most admire in
themselves, their utterance and their knowledge, is listed among the things
they were given, and the first of the gifts Chrysostom names is the remission
of sins that the Gospel of this Mass proclaims. Theophylact, following him,
repeats that grace is neither a debt nor a repayment.

Theodoret and St Thomas Aquinas keep what Chrysostom keeps, that the gifts are
God's, and add what he does not stress, that the praise is true. Theodoret
says that Paul, being about to accuse, first soothes his hearers so that the
cure may be received, and yet says nothing false, since he gives thanks to God
for gifts that had really been given them (\work{Interpretatio in I ad
Corinthios}, PG 82, col.~229). Aquinas says that the thanks are given first
\latin{ut correctionem suorum defectuum tolerabilius ferant}, so that they may
bear the correction of their faults more easily, and that the riches of
utterance and knowledge belong to the more perfect in the Church, whom the
rest share through the charity that joins them (\work{Super I ad Corinthios}
c.~1, lect.~1). The Latin commentary that goes under the name of Ambrosiaster
names the grace given in Christ as the gift by which the believer \latin{gratis
accipit remissionem peccatorum}, receives the remission of sins freely.

At the lesson's last verse the readers divide. For Chrysostom the promise that
Christ ``will confirm you unto the end without crime'' is ``also covertly
accusing them: for, to say, He shall confirm, and the word unreprovable marks
them out as still wavering, and liable to reproof''; Theodoret agrees. For
Aquinas \latin{sine crimine} means \latin{sine peccato mortali}, and the verse is
a promise of perseverance, as it is for Ambrosiaster and Cornelius a Lapide. On
either reading the confirming is the Lord's work and not the Corinthians': the
verse is a promise of it, or a reminder to the wavering that they still need
it.

\subsection{``Without thy help'': the Collect}

The Collect says outright what the Epistle implies. It asks that \latin{tuae
miserationis operatio}, the working of God's mercy, may direct our hearts, and
its reason is \latin{quia tibi sine te placere non possumus}. The two English
books that translate it, the Philadelphia missal of 1861 and the continuation
of \work{The Liturgical Year}, render \latin{sine te} ``without thy help''. The
Latin says more simply that without God himself we cannot please God. The
continuation introduces the prayer with a sentence that could stand as its
commentary: ``The surest way to obtain grace is to be ever humbly acknowledging
to our God our deep conviction that, of ourselves, we cannot please His divine
Majesty'' (vol.~XI, p.~395).

Two commentators on the Mass read the Collect in exactly this way. Berno of
Reichenau, whose Mass had another Gospel, explains it from the Communion that
follows it: \latin{Et quia Deo nec in sacrificio, nec in oblatione hostiarum
sine ejus adjutorio placere possumus, nec in atria ejus introire, nec in aula
ejus adorare eum, merito sacerdos ex sua et nostra voce Deum deprecatur,
dicens: Dirigat corda nostra} (\work{Libellus} V; PL 142, col.~1070). Without
God's help, that is, we can neither sacrifice nor enter his courts nor adore
him there; therefore the priest prays, in his own voice and ours, that God
direct our hearts. Schuster, commenting on the 1962 Mass, says that the Collect
``should establish us in humility. All the good which we do is the work of
grace, it is a gift received from God'', and he quotes St Paul, \latin{Si autem
accepisti, quid gloriaris, quasi non acceperis?} (\work{The Sacramentary} III,
p.~168). The question he quotes stands in the same letter to the Corinthians as
the Epistle (1 Cor 4:7).

\subsection{The man who was carried}

St Jerome reads the Gospel's \latin{videns Iesus fidem illorum} in its
strongest form. The paralytic was brought lying on his bed \latin{quia ipse
ingredi non valebat}, because he could not come in himself, and \latin{Videns
autem Jesus non ejus fidem qui offerebatur, sed eorum qui offerebant}: Jesus
saw not the faith of the man who was offered but the faith of those who offered
him (\work{Commentarii in Matthaeum} I, on 9:1--2; PL 26, col.~54). The reason
Jerome gives is incapacity, not merit. And he marvels at the word with which
Christ greets such a man: \latin{O mira humilitas, despectum et debilem,
totisque membrorum compagibus dissolutum, filium vocat, quem sacerdotes non
dignabantur attingere}, O wonderful humility: he calls son one despised and
feeble, loosed in every joint of his limbs, whom the priests would not deign to
touch. His moral reading follows the same line. \latin{Juxta tropologiam
interdum anima jacens in corpore suo, totis membrorum virtutibus dissolutis, a
perfecto doctore offertur curanda Domino, quae si misericordia ejus sanata
fuerit, tantum roboris accipit, ut portet statim lectulum suum}: the soul that
lies powerless in its body is offered to the Lord for healing by a perfect
teacher, and if his mercy heals it, it receives such strength that it
immediately carries its bed (col.~55). The soul is brought by another, and its
strength comes afterward.

St John Chrysostom begins from the same point and then refuses to stop there.
The bearers ``put the sick man before Christ, saying nothing, but committing the
whole to Him'', and what Christ saw, Chrysostom says, was ``the faith of them
that had let the man down. For He doth not on all occasions require faith on
the part of the sick only''. But he adds: ``Or rather, in this case the sick
man too had part in the faith; for he would not have suffered himself to be let
down, unless he had believed''. And when the scribes murmur, the one who might
have complained is silent: ``he for his part utters no such word, but gives
himself up to the power of the healer'' (\work{Homilies on Matthew} 29, 1--2).
St Thomas Aquinas, at the same verse, will not choose: \latin{Curat aliquando
Dominus aliquem propter fidem suam: aliquando propter preces suas, et aliorum},
the Lord heals sometimes for a man's own faith, sometimes for his prayers and
those of others; the bearers, read morally, are \latin{illi qui suis
monitionibus portant eum ad Deum}, those who carry the sinner to God by their
counsel (\work{Super Matthaeum} c.~9).

The disagreement between Jerome and Chrysostom is real, and Mt 9:2 does not
settle it: it says \latin{fidem illorum} and no more. Jerome's reading makes
the case at its strongest, a man saved by other people's faith, and it serves
the allegory of the nations. Chrysostom's attends more closely to the story,
since the man consented to be let down through a roof. What both readings
share, and what this interpretation rests on, is that the man did not come by
himself. Whether the faith Christ saw was the bearers' alone or his as well, it
brought him only as far as he was carried.

St Hilary, as Aquinas reports him in the \work{Catena aurea}, carries the
story over to the whole of humankind: \latin{In paralytico autem gentium
universitas offertur medenda}, in the paralytic the whole body of the nations
is offered for healing, and to it are forgiven the sins \latin{quae lex laxare
non poterat}, which the law could not loose. The nations, too, were brought.

\subsection{The one present that is not empty}

The Communion commands something: \latin{Tollite hostias, et introite}. St
Augustine names what the command asks for. ``Confession is a present unto God.
O heathen, if ye will enter into His courts, enter not empty. `Bring presents.'
What presents shall we bring with us? The sacrifice of God is a troubled
spirit: a broken and a contrite heart, `O God, shalt not Thou despise.' Enter
with an humble heart into the house of God, and thou hast entered with a
present. But if thou art proud, thou enterest empty. For whence wouldest thou be
proud, if thou wert not empty?'' (\work{Enarrationes in Psalmos} 95, 9). The
\latin{hostiae} the Communion requires are not something the worshipper owns
and hands over. They are the one thing that costs nothing but pride. Cassiodorus
reads the same word the same way, \latin{Hostias non victimas pecudum dicit,
sed conscientiae pura libamina, unde non sanguis currat, sed piae lacrymae
defluant}: not the victims of cattle but the pure offerings of conscience, from
which no blood runs but holy tears flow. He notices the order of the two
imperatives, which Augustine does not press: \latin{Sed considera quia prius
posuit, tollite, et sic introite}, first bring, and so enter, for those who do
not bring such sacrifices are not reckoned to enter the Lord's courts
(\work{Expositio psalmorum}, in Ps.~95).

\subsection{The other elements in this reading}

The Introit's word for those who pray is \latin{sustinentibus te}, those who
wait for thee: its verb is patient endurance, and the antiphon asks God to give
what waiting cannot produce. The Gradual's joy is joy at what was said to the
psalmist, \latin{in his, quae dicta sunt mihi}; in Augustine's psalter it was
joy in the persons who said it, the companions who summon the pilgrim to run.
Either way, even the gladness of the ascent is received from others. The
Alleluia's nations fear the Lord's name because God has had mercy on Sion and
built her, not because they found him: Augustine reads the verse in that order,
``Now that Thou hast pitied Sion \dots\ hence preaching hath increased among
the heathen'' (\work{Enarrationes in Psalmos} 101, s.~1, 16).

In the Offertory every finite verb has Moses for its subject. He sanctifies the
altar, offers the holocausts, immolates the victims and makes the evening
sacrifice; the people are named only at the end, \latin{in conspectu filiorum
Israel}, as those in whose sight it is done. The Church sings that antiphon over
gifts the people themselves have brought. The Secret divides its verbs in the
same way: \latin{sicut tuam cognoscimus veritatem}, as we know thy truth, is
indicative and already given; \latin{sic eam dignis moribus assequamur}, so may
we attain it by a worthy life, is subjunctive and asked for. The Postcommunion
gives thanks first and then asks to be made worthy of what has already been
received, \latin{ut dignos nos eius participatione perficias}. Its order is this
reading in small: the gift, the thanks, and then the prayer to become fit for
it. Schuster draws the practical conclusion, that a fervent Communion ``is the
best preparation for the next''.

\subsection{Not quietism: rise, take up thy bed, and walk}

Pressed too far, this reading would become quietism, and the formulary itself
refuses that. Three of its ten elements ask something of the one who prays: the
Secret's \latin{dignis moribus assequamur}, the Communion's \latin{Tollite
hostias, et introite}, and the Gospel's own \latin{Surge, tolle lectum tuum,
et vade}. Chrysostom, the principal witness of the reading, says in the same
homily that ``there is need of many labors to be able to come unto the end''
(\work{Homilies on First Corinthians} 2, 6).

The Collect itself answers the objection. It does not say that we cannot please
God; it says that we cannot please him \latin{sine te}, without him. The
distinction lies in the Latin, and it is the whole of this reading. The man in
the Gospel does rise, take up his bed and walk. What he could not do was be
forgiven, or bring himself to the house where forgiveness was given. St Thomas
Aquinas notices the exact shape of the command. The sick man had three
disabilities: he lay in bed, he was carried by others, and he could not walk;
so Christ says \latin{Surge} to the one who lay, and \latin{quia portabatur,
praecepit ut portaret}, because he was carried, he commanded him to carry
(\work{Super Matthaeum} c.~9). The one who was carried now carries what carried
him. The grace that brought him is not taken back when he stands; it is what he
stands in.

\subsection{The four senses of this reading}

\begin{fourSenses}
\item[Literal.] Paul thanks God for gifts the Corinthians did not earn and are
misusing, and promises that Christ will confirm them to the end. A paralysed
man is carried to Jesus by others and forgiven before he is healed. A psalm
tells the nations to come into the Lord's courts bringing sacrifices. The
Collect confesses that without God we cannot please him.

\item[Allegorical.] The paralytic is the Church of the nations, carried by the
faith of others into the presence of the one who forgives first and heals
after; in him the whole body of the nations is offered for healing, and its
sins, which the law could not loose, are forgiven (Hilary, as Aquinas reports
him). The one who calls such a man ``son'' shows the humility of God (Jerome).

\item[Moral.] The gifts that most tempt to pride, utterance, knowledge and
understanding, are the ones Paul lists as given, and where grace is, boasting
has no ground (Chrysostom). The soul that lies powerless is brought to the Lord
by another, and receives strength afterward (Jerome). The one present a
creature can bring is a contrite heart, and the proud enter empty (Augustine).
The one who was carried is told to carry (Aquinas).

\item[Anagogical.] What is asked for at the end is not completion but
perfection, \latin{ut dignos nos eius participatione perficias}, and the
confirming ``unto the end without crime, in the day of the coming of our Lord
Jesus Christ''. Chrysostom reads that verse as the reminder ``that not only the
beginnings must be good, but the end also \dots\ Therefore there is need of
patience''.
\end{fourSenses}
